\scanpage{9}
\mspage{2}
% margin note written slanted, upper left
\margin{Voir p. \unsure{90|9} \unsure{1bis}, 2bis, à mettre \unsure{avant|au v.} \unsure{les notes|les suites|leur suite} sur les figures.}

On appelle \uline{figure} \inserted{(inconditionnelle)} d'un magasin, un élément de l'ens
\[
\unsure{\mathcal F|\mathscr F}_M=\{\unsure{\Phi|\varphi}\subset M\mid\Phi\ \text{fermé dans $M$ pour}\ \leq,\ \text{et}\ \forall X,Y\in\Phi,\ \text{on a}\ X\mathrel{\unsure{\asymp|\bowtie|\gtrless}}Y\}
\]
L'ens $\mathcal F_M$ est l'ensemble des ``figures inconditionnelles'' (ou ``figures'' tout court) du magasin.

% margin note written slanted, left of the preceding paragraph
\margin{$\mathcal F_M$ sera appelé \uline{\unsure{Nablien|Nablium|Nablum}} \uline{\unsure{universel}} associé au magasin $M$.}

On sera amené, à l'occasion, de faire sur les figures envisagées deux types possibles de restrictions

1°) Conditions du type ``finitude'' ou ``locale finitude''.

\deleted{\unsure{sur les}}

2°) Condition restrictive ``polyédrale'' : si $X,Y\in\Phi$ et $\widetilde X\cap\widetilde Y\ne\varnothing$, on veut qu'il \deleted{$\widetilde X\cap\widetilde Y$} admette un plus grand élément, i.e. soit de la forme $\widetilde Z$, ou encore que $\operatorname{Inf}(X,Y)$ existe.

On étend à $\mathcal F$ les relations d'ordre $\leq$ et $\ll$ de $M$ (en plongeant $M$ dans $\mathcal F$ via les ``figures élémentaires'' $\widetilde X$, $X\in M$
\diagram{\begin{tikzcd}[ampersand replacement=\&]
M\arrow[r,hook]\&\mathcal F_M
\end{tikzcd}}
\[
\deleted{\unsure{\widetilde{\vphantom X}}}X\mathrel{\unsure{\mapstochar\relbar\relbar\bullet|\longmapsto}}\widetilde X=M_{\leq X}\qquad-
\]
\unsure{souvent} et par abus de langage, on identifie $M$ à \unsure{l'ens} des figures élémentaires, $X\in M$ : la figure élémentaire associée...)

% margin note written slanted, down the lower left edge
\margin{\unsure{Dans ce qui suit} on \unsure{désidentifiera} les figures d'avec les parties de $M$ qui les définissent, et pour une figure $F$ \unsure{désormais}, on \unsure{identifiera|identifie|désignera} la partie correspondante $\widetilde F=\{X\in M(\unsure{\subset|\in}\mathcal F)\mid X\leq F\}$.}

Donc
\[
F\leq G\overset{\mathrm{def}}{\Longleftrightarrow}F\subset G
\]
\[
\begin{aligned}
\text{NB le Sup des $F$ pour $\leq$ est}\quad&\bigcup\quad\text{dans $\unsure{P|\mathcal P|\mathscr P}(M)$}\\
\text{le Inf}\quad\underline{\hspace{3.3cm}}\quad&\bigcap\quad\underline{\hspace{2cm}}
\end{aligned}
\]
\[
F\ll G\overset{\mathrm{def}}{\Longleftrightarrow}\forall\unsure{x|X}\in\widetilde F,\ \exists Y\in\widetilde G,\ \text{avec}\ X\ll Y
\]
\[
\begin{aligned}
F\asymp G&\Longleftrightarrow\forall X\in\widetilde F,\ Y\in\widetilde G,\ \text{on a :}\ X\asymp Y\\
&\Longleftrightarrow\{F,G\}\ \text{majoré dans}\ \mathcal F\\
&\Longleftrightarrow\operatorname{Sup}(F,G)\ \text{existe dans}\ \mathcal F\ \deleted{\unsure{\text{(et }F\cup G\text{)}}}\\
&\Longleftrightarrow F\cup G\in\mathcal F
\end{aligned}
\]
% archive sheet number: 9
